\chapter{Style, theme, and the palette}
\label{ch:style}

\begin{aside}
Colours and fonts are the bit users see first. A library that
treats them carelessly produces UIs that look wrong on half the
terminals they run on. \maya{}'s style system is more elaborate
than the typing-friendly DSL might suggest: a palette of named
roles, a theme that maps roles to colours, a style struct that
overrides per-element, and a merger that combines them under a
defined precedence.
\end{aside}

\section{Three layers of style}

When the renderer paints a cell, the cell's style is the result
of three layers, applied in order:

\begin{enumerate}[leftmargin=*]
  \item \textbf{Theme.} The application-level palette. Maps named
    roles (\emph{foreground}, \emph{accent}, \emph{warning})
    to concrete colours.
  \item \textbf{Inherited style.} The style applied by parent
    boxes via the DSL pipe or runtime equivalent.
  \item \textbf{Local style.} The style attached to the
    text/box at the leaf.
\end{enumerate}

Each layer can override the previous. The merge rule
(Chapter~\ref{ch:elements}) decides per-attribute whether to
inherit or override.

\subsection*{Why three layers?}

Two layers (theme + override) would work for simple cases. The
third (inheritance) is what makes ``the whole panel is bold''
work without setting Bold on every leaf.

The trade-off: more layers means more bookkeeping. \maya{}'s
merge is a small constant-time operation per cell, so the cost
is acceptable.

\section{The palette}

A palette is a fixed-size array of \code{Color} values, each
indexed by a named role. \maya{}'s default palette has 24 slots:

\begin{itemize}[leftmargin=*]
  \item \code{Foreground}, \code{Background}: the default
    foreground and background.
  \item \code{Muted}: dim text (timestamps, labels).
  \item \code{Subtle}: secondary text (placeholders).
  \item \code{Primary}: the brand colour.
  \item \code{Accent}: highlight colour.
  \item \code{Success}, \code{Warning}, \code{Danger},
    \code{Info}: semantic colours.
  \item \code{Border}: default border colour.
  \item \code{Selection}, \code{Cursor}: text-input UI.
  \item Ten ``user'' slots for application-specific roles.
\end{itemize}

A palette is just a \code{std::array<Color, 24>}, plus a small
enum for the indices. Looking up is constant time.

\subsection*{Roles vs colours}

The point of the palette is to decouple ``what the cell means''
from ``what colour to draw''. Code that wants a warning highlight
asks for \code{Palette::Warning}; the theme decides whether that
is amber, orange, or yellow.

This is the same indirection CSS uses (via custom properties) and
that design systems (Tailwind, Material Design) build whole
ecosystems around.

\section{The Theme struct}

A theme is a palette plus metadata:

\begin{lstlisting}
struct Theme {
    Palette  palette;
    std::string name;        // "dark", "light", "solarized"
    bool     is_dark;        // light-mode UIs flip choices
    Border   default_border; // BorderStyle::Rounded by default
};
\end{lstlisting}

\maya{} ships three themes out of the box: \code{Dark} (default),
\code{Light}, and \code{HighContrast}. Loading a theme replaces
the palette wholesale.

\subsection*{Defining a theme}

\begin{lstlisting}
constexpr Theme my_theme = {
    .palette = {
        [Role::Foreground] = Color::rgb(220, 220, 240),
        [Role::Background] = Color::default_(),
        [Role::Muted]      = Color::rgb(120, 130, 150),
        [Role::Primary]    = Color::rgb(120, 180, 250),
        // ...
    },
    .name = "midnight",
    .is_dark = true,
};
\end{lstlisting}

Designated initialisers (C++20) let you skip slots you do not
care about; they retain their default values.

\section{The Style struct, revisited}

\begin{lstlisting}
struct Style {
    Color fg = Color::default_();
    Color bg = Color::default_();
    Attribute attributes = Attribute::None;
    std::optional<std::string> hyperlink;
};
\end{lstlisting}

We saw this in Chapter~\ref{ch:elements}. The colour fields can
be a literal RGB, an Ansi16 index, an Ansi256 index, or
``default'' (let the terminal pick).

A style created from a role:

\begin{lstlisting}
Style accent_style = Style::from_role(Role::Accent);
\end{lstlisting}

This is a thin helper that looks up the colour in the current
theme. The lookup is deferred until paint --- if the theme
changes between frames, the style picks up the new value.

\subsection*{Attribute, in detail}

\code{Attribute} is a bitset of typographic attributes:

\begin{itemize}[leftmargin=*]
  \item \code{Bold}: SGR 1.
  \item \code{Dim}: SGR 2.
  \item \code{Italic}: SGR 3.
  \item \code{Underline}: SGR 4.
  \item \code{Blink}: SGR 5. (Rarely supported well; avoid.)
  \item \code{Reverse}: SGR 7. Swaps fg and bg.
  \item \code{Strike}: SGR 9.
  \item \code{DoubleUnderline}: SGR 21.
  \item \code{Overline}: SGR 53. (Even more rarely supported.)
\end{itemize}

Stored as a small bitmask. Two attributes combine with bitwise
OR; this is what makes \code{Bold | Italic} mean ``both''.

\section{Style merge: the precedence rule}

The merge rule, again, with the formal precedence:

\begin{enumerate}[leftmargin=*]
  \item \textbf{Attributes.} Result is the union of both. Bold
    in either input means Bold in the output.
  \item \textbf{Colours.} If the inner style sets a non-default
    colour, the inner wins. Otherwise, the outer's colour
    propagates.
  \item \textbf{Hyperlink.} Inner wins if set.
\end{enumerate}

In code:

\begin{lstlisting}
Style merge(const Style& outer, const Style& inner) {
    Style r = outer;
    r.attributes |= inner.attributes;
    if (!inner.fg.is_default()) r.fg = inner.fg;
    if (!inner.bg.is_default()) r.bg = inner.bg;
    if (inner.hyperlink)        r.hyperlink = inner.hyperlink;
    return r;
}
\end{lstlisting}

The asymmetry (attribute union, colour override) is what makes
the system intuitive. ``Inherit the Bold from the parent and the
colour from the child'' is what most users mean.

\begin{idea}
This is similar to CSS cascade rules, simplified. CSS has six
levels of specificity; \maya{} has two (outer, inner) and one
rule (attribute-union, colour-override). The simpler rule wins
in clarity.
\end{idea}

\section{Themes at runtime}

To change theme:

\begin{lstlisting}
maya::set_theme(Theme::Light());  // or any custom theme
\end{lstlisting}

The pipeline notices on the next frame and reflows. Any style
that referenced a role updates automatically.

\subsection*{Persisting theme preferences}

If you store the theme name in a config file, you can restore it
on startup:

\begin{lstlisting}
auto config = load_config();
maya::set_theme(themes::by_name(config.theme).value_or(Theme::Dark()));
\end{lstlisting}

The pipeline does not care where the theme came from; it just
applies whatever is current.

\section{High-contrast and accessibility}

The \code{HighContrast} theme is designed for users who need
sharp distinctions: pure black background, pure white foreground,
saturated accent colours. The palette is calibrated to a minimum
WCAG-AA-equivalent contrast ratio.

If your app respects \code{NO\_COLOR=1} (an environment variable
that means ``do not use colour''), \maya{} will use a theme that
relies on attributes (Bold, Underline) instead of colour:

\begin{lstlisting}
if (std::getenv("NO_COLOR")) {
    maya::set_theme(themes::Monochrome());
}
\end{lstlisting}

\begin{aside}
The NO\_COLOR convention was proposed by Jason Boice in 2017 and
is now widely supported (ls, ripgrep, fd, bat, hyperfine,
\maya{}). When a user sets it, they mean it: drop all colour.
\end{aside}

\section{Detecting terminal capabilities}

Not every terminal supports truecolour, or italics, or OSC 8.
\maya{} detects capabilities at startup:

\begin{itemize}[leftmargin=*]
  \item \code{COLORTERM=truecolor} or \code{=24bit}: emit
    full RGB.
  \item \code{TERM} starts with one of a known set (xterm-256,
    \emph{etc.}): use 256-colour palette.
  \item Otherwise: 8/16-colour palette.
\end{itemize}

The detection happens once in
\file{maya/include/maya/terminal/terminal.hpp} and the result
threads through the pipeline as a \code{TerminalCaps} struct.

When the renderer emits a colour, it consults caps:

\begin{itemize}[leftmargin=*]
  \item RGB asked but only 256 supported: quantise.
  \item RGB asked but only 16 supported: nearest-neighbour to
    16-colour palette.
  \item Default-fg asked: emit nothing (terminal's choice).
\end{itemize}

\section{Building your own theme}

If you want to ship a theme with your app:

\begin{lstlisting}
namespace mythemes {
    inline constexpr Theme Sunset = {
        .palette = {
            [Role::Foreground] = Color::rgb(245, 230, 200),
            [Role::Background] = Color::rgb(40, 25, 35),
            [Role::Primary]    = Color::rgb(220, 140, 80),
            [Role::Accent]     = Color::rgb(255, 180, 120),
            [Role::Warning]    = Color::rgb(255, 200, 100),
            [Role::Danger]     = Color::rgb(220, 80, 80),
            [Role::Success]    = Color::rgb(140, 200, 130),
            [Role::Muted]      = Color::rgb(160, 140, 130),
            [Role::Border]     = Color::rgb(120, 90, 80),
        },
        .name = "sunset",
        .is_dark = true,
    };
}
\end{lstlisting}

Then set it:

\begin{lstlisting}
maya::set_theme(mythemes::Sunset);
\end{lstlisting}

Every styled element using \code{Role::Primary} will use
RGB(220, 140, 80). No code change required.

\section{Style preset gallery}

Reusable style preset patterns:

\begin{lstlisting}
constexpr auto H1 = Bold | Underline | Fg<255, 240, 200>;
constexpr auto H2 = Bold | Fg<200, 220, 240>;
constexpr auto Body = Style{};
constexpr auto Code = Italic | Fg<200, 200, 160>;
constexpr auto Subtle = Dim | Fg<140, 140, 150>;
constexpr auto Error = Bold | Fg<255, 100, 100>;
constexpr auto Success = Bold | Fg<100, 200, 130>;
\end{lstlisting}

Each is a \code{StyTag} at the type level. Reusing them keeps
your app's appearance consistent and lets you tweak in one
place.

\section{Pitfalls}

\begin{warning}
\textbf{Hard-coding colours instead of using roles.} If you write
\code{Fg<255, 200, 100>} everywhere, you cannot theme the app.
Use \code{role::Warning} (or the equivalent helper) so a theme
swap actually changes the appearance.
\end{warning}

\begin{warning}
\textbf{Setting BG to black on a non-default terminal.} Some
users have dark grey or off-black backgrounds. Setting your bg
to RGB(0,0,0) creates a visual line. Use \code{Color::default\_()}
unless you really need a specific bg.
\end{warning}

\begin{warning}
\textbf{Forgetting NO\_COLOR.} If your app relies on colour to
convey meaning (red = bad, green = good), users with
\code{NO\_COLOR=1} will lose the cue. Pair colour with text
or attributes (``[ERR]'', Underline).
\end{warning}

\begin{warning}
\textbf{Truecolour quantisation surprises.} A theme that uses
adjacent RGB values may collapse onto the same Ansi256 cell on
old terminals. Test your theme on a 256-colour terminal before
shipping.
\end{warning}

\section*{Workshop}

\begin{enumerate}[leftmargin=*]
  \item Define a \code{Theme} of your own with three named roles.
    Use it with \code{maya::set\_theme}. Toggle between yours
    and the default at runtime with a hotkey.
  \item Inspect what your terminal advertises: print
    \code{getenv("TERM")} and \code{getenv("COLORTERM")} in a
    counter app.
  \item Run with \code{NO\_COLOR=1} set and confirm the app
    still conveys distinctions through attributes.
\end{enumerate}

\begin{recap}
The style system has three layers (theme, inheritance, local) and
merges with a defined precedence (attribute-union,
colour-override). The palette decouples roles from colours;
themes swap palettes wholesale. Truecolour quantisation,
NO\_COLOR, and capability detection happen at the edges so your
code can think in roles, not raw RGB.
\end{recap}
